\chapter{系统基础}

\section{领域导言与演进背景}

\begin{preamblebox}[知识领域概述与变革]
计算机系统是由底层硬件与上层系统软件共同构建的综合计算基础设施，是所有现代应用软件运行的物理与逻辑载体。在当今数字化社会中，计算机系统已成为支撑人类生产与生活不可或缺的中枢基石。因此，系统学习计算机系统全栈知识、熟练掌握系统设计与工程构建技能，并深刻洞察其底层运行原理与核心设计哲学，是培养计算机专业拔尖人才的核心要务。

在计算机科学专业课程体系中，系统知识的传授通常跨越操作系统、并行与分布式系统、计算机网络、计算机体系结构与组织以及软件工程等多门核心骨干课程。正如其名称所示，“系统基础 (SF)”知识领域旨在提炼这些系统类核心课程共有且通用的基础概念与思维模式，以高度凝练、前后贯通的统一视角建立宏观全景，为学生后续深入学习各专用领域的特定机制与策略奠定坚实的系统思维基石。

\textbf{自 CS2013 以来的核心变革：}
\begin{enumerate}
    \item \textbf{增设两大全新核心单元}：正式增设“系统安全”与“系统设计”两个独立的知识单元，将现代安全防御与宏观架构设计准则确立为必修基石。
    \item \textbf{重组并创立“系统性能”单元}：废弃原 CS2013 中的“临近性”与“虚拟化与隔离”单元，将其核心精髓整合提升为全新的“系统性能”单元。
    \item \textbf{重组并创立“性能评估”单元}：废弃原 CS2013 中的“评估”与“定量评估”单元，全面重构为现代化的“性能评估”单元。
    \item \textbf{计算范式演进为“计算机系统概述”}：将原“计算范式”单元调整更名为“计算机系统概述”，精简陈旧内容，并吸收并入了原“跨层通信”单元的核心主题。
    \item \textbf{更名并深化“系统基本概念”}：将原“状态与状态转换”更名为“系统基本概念”，系统性补充了有限状态机（FSM）与时序系统建模等理论内容。
    \item \textbf{优化调整并行性主题}：精简了原“并行性”单元中的部分内容，重点保留并强化了简易应用级并行处理（如请求级并发、多线程多服务器模型及流水线）。
    \item \textbf{废弃独立跨层通信单元}：将原“跨层通信”相关主题有机并入“计算机系统概述”单元。
    \item \textbf{废弃原双重评估单元}：将原“评估”与“定量评估”合并提炼至“性能评估”单元。
    \item \textbf{整合临近性与虚拟化主题}：将原“临近性”与“虚拟化与隔离”的议题无缝融入“系统性能”单元。
    \item \textbf{废弃独立并行单元并归并基础概念}：取消独立的“并行性”单元，将其基础部分融入“系统基本概念”单元。
    \item \textbf{重命名并升格“系统可靠性”}：将原“基于冗余的可靠性”重命名并拓展为具有全生命周期视角的“系统可靠性”单元。
\end{enumerate}
\end{preamblebox}

\section{核心学时分配建议}

本知识领域在“向日葵模型”下的核心学时推荐如下表所示（计算机核心共 18 学时，领域核心共 9 学时，总计 27 学时）：

\begin{table}[H]
\centering
\setlength{\tabcolsep}{3.5pt}\fontsize{8pt}{9.6pt}\selectfont\begin{tabularx}{\columnwidth}{Y c c}
\toprule
\textbf{知识单元} & \textbf{计算机核心} & \textbf{领域核心} \\
\midrule
计算机系统概述 & 3 & -- \\
系统基本概念 & 4 & -- \\
资源分配与调度 & 1 & 2 \\
系统性能 & 2 & 2 \\
性能评估 & 2 & 2 \\
系统可靠性 & 2 & 1 \\
系统安全 & 2 & 1 \\
系统设计 & 2 & 1 \\
\midrule
\textbf{合计} & \textbf{18 学时} & \textbf{9 学时} \\
\bottomrule
\end{tabularx}
\end{table}

\section{知识单元与详细大纲}

\subsection{计算机系统概述}
\textbf{【学时建议】}：3 计算机核心学时。

本单元建立对现代计算系统分层架构与软硬件协同设计的宏观全局认知，阐释从基础物理逻辑到高层应用程序接口的逐级抽象演进。

\subsubsection{核心主题}
\begin{itemize}
    \item \textbf{计算机基本构建块与组件}：逻辑门电路、触发器、寄存器、总线与内部互连结构；现代处理器的经典划分（数据通路 Datapath + 控制逻辑 Control + 存储体系 Memory）。
    \item \textbf{作为计算范式的底层硬件}：基本数字逻辑单元；布尔代数逻辑表达式、卡诺图化简、积之和式（与或式，Sum-of-Products）。
    \item \textbf{程序设计抽象与接口规范}：软件抽象层级划分、应用程序编程接口（API）、底层系统库与动态链接库调用机制。
    \item \textbf{应用服务与操作系统服务的边界辨析}：用户态与内核态特权级分离、内核系统服务职责边界、远程过程调用（RPC）核心思想。
    \item \textbf{应用层与操作系统交互机制}：系统调用时序、陷阱指令、上下文保存与恢复、跨特权级数据拷贝开销。
    \item \textbf{流水线技术基础原理}：指令与任务级流水线核心概念、重叠执行阶段划分与吞吐提升。
    \item \textbf{系统扩展性基本概念}：强扩展（Strong Scaling，算力扩充加速单任务求解）与弱扩展（Weak Scaling，资源扩充支持更大问题规模）的本质差异。
\end{itemize}

\subsubsection{示范性学习成果与能力目标}
\begin{enumerate}
    \item \textbf{[解释]} 详细阐明计算机的基本硬件构建模块及其在计算机体系结构历史演进过程中的关键作用。
    \item \textbf{[应用/开发]} 运用基本逻辑门与触发器等硬件设计基础构建块，设计并实现可解决基础计算任务的组合与时序逻辑电路（例如全加器、计数器）。
    \item \textbf{[评估]} 熟练运用原理图捕获、逻辑综合与数字仿真工具链，对数字逻辑电路设计的逻辑功能与时序特性进行验证与性能评估。
    \item \textbf{[解释]} 深刻阐释现代计算系统如何基于“关注点分离”设计哲学构建层级体系，通过严谨规范的接口协议向上层屏蔽底层实现的复杂细节。
    \item \textbf{[解释]} 深入阐明底层硬件、操作系统内核、虚拟机运行时与应用层软件在计算本质上均是提供额外解释与加工能力的多层递进处理实体。
    \item \textbf{[解释]} 阐明在复杂分层体系结构中，底层故障检测、跨层异常状态向上汇报以及分级差错恢复与容错降级的运行机制。
    \item \textbf{[开发]} 综合运用分层设计方法、差错检测与自愈恢复技术，编写实现一个能够跨层反馈并优雅处理错误状态的典型软件系统（如可靠的 TCP 客户端/服务器通信系统）。
    \item \textbf{[应用/评估]} 熟练运用现代程序跟踪分析、单步调试及断点分析工具，在具有深层调用栈的分层系统中精准定位并修复复杂软件缺陷。
    \item \textbf{[解释/评估]} 透彻理解强扩展与弱扩展的本质机理，即问题求解性能如何受制于问题规模与计算资源配比，并能结合现实云计算或超级计算工程实例进行深度阐述。
\end{enumerate}

\subsection{系统基本概念}
\textbf{【学时建议】}：4 计算机核心学时。

本单元提炼计算系统中最本质的数学与物理抽象，涵盖连续与离散、状态与状态机、时钟时序以及串行与并行计算范式。

\subsubsection{核心主题}
\begin{itemize}
    \item \textbf{系统范式对比}：数字系统与模拟系统、离散系统与连续系统的本质差异与现实工业对应。
    \item \textbf{布尔逻辑与化简}：基本逻辑门、复合逻辑表达式求值、布尔代数恒等律与真值表化简。
    \item \textbf{时钟、状态与时序控制}：系统时钟信号、触发边沿、时钟周期节拍、状态保持与时序同步机制。
    \item \textbf{状态与状态转移}：离散系统状态空间定义、初始状态、中间态、终止状态以及状态全生命周期演化模型。
    \item \textbf{有限状态机理论}：确定性有限状态自动机（DFA）与非确定性有限状态自动机（NFA）的形式化定义、状态转移表与状态转移图。
    \item \textbf{逻辑电路与存储体系}：无记忆组合逻辑电路与含反馈记忆时序逻辑电路的辨析；锁存器、寄存器、静态与动态存储器单元。
    \item \textbf{作为状态机的计算机与协议}：通用计算机体系结构作为指令解释状态机；网络通信协议（如 TCP 三次握手）作为协同状态机。
    \item \textbf{串行与并行处理范式辨析}：顺序计算流与并发计算流的本质区别。
    \item \textbf{应用级串行处理}：单线程执行模型、程序计数器顺序推进与阻塞式 I/O。
    \item \textbf{简易应用级并行处理}：
    \begin{itemize}
        \item 请求级并行（Web 服务、C/S 架构客户端-服务器端交互、分布式集群服务）。
        \item 单服务器单线程轮询事件循环模型。
        \item 多服务器多线程并发处理架构。
        \item 任务级流水线并行及其重叠加速机制。
    \end{itemize}
\end{itemize}

\subsubsection{示范性学习成果与能力目标}
\begin{enumerate}
    \item \textbf{[解释]} 深刻辨析数字系统与模拟系统、离散系统与连续系统的根本区别，并能列举工程实际与日常生活中的典型实例进行说明。
    \item \textbf{[解释]} 将通用的计算过程形式化建模为具备已知配置集合（状态空间）并在不同配置间依照确定规则发生跃迁（状态转移）的动态演化系统。
    \item \textbf{[解释]} 严密剖析输出仅取决于当前输入的无状态系统与具有历史记忆依赖的有状态系统在架构设计、并发控制与横向扩展上的本质异同。
    \item \textbf{[开发]} 针对具体工程控制问题（如多相位交通信号灯时序控制器、通信数据流序列模式匹配器），建立形式化的有限状态机状态转移模型与规范。
    \item \textbf{[解释]} 将通用计算机体系结构深刻解释为一个逐条取指、译码并依据操作码更新内部寄存器与内存状态的确定性状态机。
    \item \textbf{[解释]} 阐明通用计算机软件程序或网络通信协议如何通过状态机进行精确无歧义建模，并论证同一计算逻辑可存在多种等价的状态表示方式。
    \item \textbf{[分析/评估]} 根据给定的状态机模型规范（如 TCP 连接生命周期管理状态机），严密推导并绘制系统在各种正常与异常输入场景下的时序演变行为曲线。
    \item \textbf{[开发]} 针对同一计算问题，分别编写其单线程串行版本和简单多线程/多进程并行版本的程序实现。
    \item \textbf{[评估]} 在不同输入规模负载下，评测并对比串行程序与并行程序的运行耗时，准确计算并量化其实际达成的并行加速比。
    \item \textbf{[解释/评估]} 在执行时间轴图上可视化展示并发操作与并行事件在物理时间上的同时发生，清晰阐述发掘并发潜力如何有效压缩任务的端到端运行耗时。
\end{enumerate}

\subsection{资源分配与调度}
\textbf{【学时建议】}：1 计算机核心学时，2 领域核心学时。

本单元深入剖析计算机有限物理与虚拟资源的分配机制，研究兼顾吞吐量、低延迟、公平性与防饥饿的经典调度算法及其工程权衡。

\subsubsection{核心主题}
\begin{itemize}
    \item \textbf{多维度计算机资源分类}：CPU 处理器时间配额、物理/虚拟内存空间、磁盘/固态硬盘存储 I/O 带宽、网络传输带宽及专用硬件加速器。
    \item \textbf{经典调度算法全景}：
    \begin{itemize}
        \item 先来先服务调度。
        \item 基于优先级的调度（静态优先级、动态优先级调整）。
        \item 最早截止时间优先调度（EDF，实时系统）。
        \item 公平份额调度与完全公平调度器（CFS）。
        \item 抢占式调度与非抢占式调度机制。
    \end{itemize}
    \item \textbf{常见调度算法的优缺点权衡}：周转时间、系统吞吐率、首响响应时间、调度决策开销、算法复杂度、资源利用率以及线程饥饿和优先级反转的防范。
\end{itemize}

\subsubsection{示范性学习成果与能力目标}
\begin{enumerate}
    \item \textbf{[解释]} 深刻阐释有限的计算机软硬件资源（CPU 份额、物理内存、持久化存储以及网络带宽）如何通过细粒度配额分配与准入控制有效供给多个竞争实体共享。
    \item \textbf{[解释]} 详细阐述主流调度算法（FCFS、轮转 RR、抢占式优先级调度、公平分享调度）的核心数据结构、工作原理与调度时序推进机制。
    \item \textbf{[评估]} 全面分析并评估常见调度算法在不同工作负载特征下的优劣势，深入权衡吞吐量、响应延迟、调度器开销与公平性之间的设计取舍。
    \item \textbf{[开发/评估]} 编程实现典型调度算法原型，搭建合成测试工作负载环境，对其调度延迟、吞吐指标与公平性表现进行严密的量化评测与对比。
\end{enumerate}

\subsection{系统性能}
\textbf{【学时建议】}：2 计算机核心学时，2 领域核心学时。\\
\textbf{【跨领域交叉引用】}：体系结构/Memory Management（内存管理）, 操作系统/Virtual Memory（虚拟内存）。

本单元从软硬件结合视角剖析物理时延鸿沟与局部性原理，揭示层次存储、多级缓存以及虚拟化与隔离机制对系统性能的关键赋能。

\subsubsection{核心主题}
\begin{itemize}
    \item \textbf{计算机系统中的物理与逻辑时延}：
    \begin{itemize}
        \item 物理光速极限与纳秒级计算尺度：真空光速约为每纳秒 30 厘米（1 英尺），相比之下 1~GHz 处理器的主频周期即为 1 纳秒。
        \item 数量级访存鸿沟对比：CPU 内部寄存器（$<1$~ns）、L1/L2/L3 高速缓存（$1\sim 10$~ns）、主存 DRAM 访问（$50\sim 100$~ns）、固态硬盘与机械磁盘 I/O（微秒至毫秒级）、跨局域网与广域网远程网络通信（毫秒至百毫秒级）。
    \end{itemize}
    \item \textbf{多级缓存机制与局部性原理}：时间局部性与空间局部性对处理器执行流水线及全栈系统吞吐的决定性影响。
    \item \textbf{全栈缓存与缓存一致性}：
    \begin{itemize}
        \item 缓存抽象的普遍性：数据库查询缓存、操作系统虚拟文件系统页缓存、分布式键值缓存以及多核体系结构硬件缓存。
        \item 缓存一致性挑战与维护机制（写穿透、写回、总线监听与目录式一致性协议）。
    \end{itemize}
    \item \textbf{分层存储系统与平均访存时间}：处理器分层存储体系结构引论；平均访存时间经典计算公式：$\text{AMAT} = \text{Hit Time} + \text{Miss Rate} \times \text{Miss Penalty}$。
    \item \textbf{虚拟化与隔离的底层设计动因}：多租户运行下的安全边界防护、单应用故障隔离以及可预测的服务质量保障。
    \item \textbf{间接层 设计哲学}：“计算机科学领域的任何问题都可以通过增加一个间接层来解决”；以虚拟内存对物理内存的解耦与重映射管理为例深入分析。
    \item \textbf{虚拟内存与虚拟机的实现机制}：页表、多级页表、转换后援缓冲器（TLB 快表）、缺页异常处理机制；虚拟机监控器、硬件辅助虚拟化技术（如 Intel VT-x）与容器化轻量级隔离技术。
\end{itemize}

\subsubsection{示范性学习成果与能力目标}
\begin{enumerate}
    \item \textbf{[解释]} 深刻阐释时间局部性与空间局部性原理在决定计算机底层硬件微架构、算法运行效率与系统软件性能中的基石作用。
    \item \textbf{[计算/评估]} 熟练运用平均访存时间（AMAT）公式进行严密量化推导，从存储容量、命中率/缺失率、访问时延与硅片开销四个维度深度评估存储层次结构的性能权衡。
    \item \textbf{[解释]} 深入阐明在多个并发程序与互不信任的执行环境共享底层物理计算资源的场景下，实施强隔离与安全保护的核心必要性。
    \item \textbf{[解释]} 详细阐述间接重映射抽象如何为并发运行的多道程序与隔离环境创造出拥有“专属独占物理机器与连续内存资源”的透明逻辑幻象。
    \item \textbf{[评估]} 针对运行在相互隔离的独立虚拟机或容器环境中的两个应用实例，实测并量化其在共用 CPU/内存/磁盘高压负载下的相互干扰度，科学评估系统性能隔离机制的有效性。
\end{enumerate}

\subsection{性能评估与评测}
\textbf{【学时建议】}：2 计算机核心学时，2 领域核心学时。

本单元建立现代计算机系统的科学量化评估方法学，涵盖核心指标体系构建、基准测试集设计、性能分析定律、数量级估算以及性能诊断工具链。

\subsubsection{核心主题}
\begin{itemize}
    \item \textbf{系统性能核心度量指标体系}：系统吞吐量、端到端响应时间、长尾延迟（Tail Latency, P99/P99.9 分位数）、资源饱和度与能效比（吞吐/瓦特）。
    \item \textbf{工作负载建模与代表性基准测试集}：标准合成负载建模、代表性工业基准测试套件（如 SPEC CPU、TPC-C、YCSB 等），以及性能测量数据的规范化采集、统计滤波与分布特征分析。
    \item \textbf{CPI (Cycles per Instruction) 性能方程}：$\text{CPU Time} = I \times \text{CPI} \times T_{\text{clk}}$；作为深刻理解指令集架构演进、处理器流水线深度与层次存储系统组织权衡的定量分析工具。
    \item \textbf{阿姆达尔定律 (Amdahl's Law)}：计算任务中无法被加速或无法并行的串行瓶颈部分，严格限制了增加算力资源所能获得的整体最大加速比上限：$S(s) = \frac{1}{(1-p) + \frac{p}{s}}$。
    \item \textbf{指导定量评估的解析建模工具}：排队论模型直觉（Little 定律、M/M/1 队列模型）、均值分析与系统瓶颈理论推演。
    \item \textbf{数量级分析与大 O 渐进记号}：Order of Magnitude 分析法、大 O 渐进时间/空间复杂度分析及其在工业级大规模系统资源预算规划中的应用。
    \item \textbf{快路径与慢路径分离}：系统架构中的“快路径”与“慢路径”设计哲学（让常见关键路径极度精炼高效，慢路径保障完备性）。
    \item \textbf{底层系统微架构事件对宏观性能的损耗}：流水线冒险停顿、多级缓存缺失、缺页异常、TLB 缺失刷新与分支预测错误惩罚。
    \item \textbf{分层系统与多租户异构平台评测挑战}：复杂分层调用栈、并发多租户工作负载与异构计算平台的非线性相互作用，及其给客观定量评测带来的工程挑战。
    \item \textbf{微基准测试 的常见陷阱}：编译器激进优化剔除关键测试代码、系统冷启动与缓存预热偏差、操作系统动态调频（DVFS）干扰与多核伪共享（False Sharing）。
\end{itemize}

\subsubsection{示范性学习成果与能力目标}
\begin{enumerate}
    \item \textbf{[解释]} 深入阐明计算机系统体系结构各层组件（硬件微架构、操作系统内核调度、语言运行时、网络栈）如何协同促进系统端到端性能的提升。
    \item \textbf{[解释]} 阐明在特定工业场景（如高频量化交易系统、大规模交互式 Web 搜索、大数据离线批处理计算）下选用特定性能指标（如平均时延 vs 尾延迟 P99、批处理总吞吐 vs 单任务执行时间）的判定依据与业务意义。
    \item \textbf{[解释/评估]} 辩证阐释基准测试集在客观评价系统性能方面的核心指导价值及其固有的代表性偏差与过度优化风险。
    \item \textbf{[解释/分析]} 准确阐述阿姆达尔定律的数学公式与物理内涵，深入讨论其在强扩展与大规模异构并行计算环境下的适用边界与工程局限。
    \item \textbf{[应用/评估]} 熟练运用极限分析法或粗略数量级心算法，在复杂系统方案规划阶段快速定量估算出目标系统性能指标的理论上下界。
    \item \textbf{[应用]} 熟练运用现代系统级性能剖析与诊断工具链（如 perf、eBPF、火焰图 FlameGraph、gprof 等），对目标系统的执行热点、调用链路与资源消耗进行精细化测量与画像。
    \item \textbf{[开发/评估]} 针对典型主流工业级计算系统（如 Linux 操作系统内核、分布式大数据分析引擎 Spark），自主设计并端到端执行一套规范严密的性能基准评测实验。
    \item \textbf{[评估]} 在复杂分层计算系统环境中开展受控对照基准实验，精确量化某一系统底层关键参数（如并发线程池规模、缓存淘汰策略、TCP 拥塞算法）调整对整体宏观性能产生的实际影响。
\end{enumerate}

\subsection{系统可靠性}
\textbf{【学时建议】}：2 计算机核心学时，1 领域核心学时。

本单元立足全生命周期的系统高可用与容错工程，辨析缺陷、差错与故障的级联传导，掌握基于校验、冗余编码与状态复制的容错机制。

\subsubsection{核心主题}
\begin{itemize}
    \item \textbf{软件缺陷与物理故障辨析}：程序逻辑缺陷、系统运行时差错与底层硬件物理故障（Faults，如坏内存位翻转、瞬态电磁干扰）的本质机理与级联演变路径。
    \item \textbf{基于冗余保障高可靠性机制之一：状态校验与重试}：校验和、超时与自适应重传、状态回滚与幂等操作设计。
    \item \textbf{基于冗余保障高可靠性机制之二：信息冗余编码}：检错与纠错编码（海明码 ECC、循环冗余校验 CRC、前向纠错码 FEC）。
    \item \textbf{基于冗余保障高可靠性机制之三：空间与硬件冗余}：双机热备、磁盘阵列镜像容错（RAID-1/RAID-5/6）、多副本状态机复制。
    \item \textbf{系统可靠性与高可用其他工程途径}：健康心跳探测、看门狗定时器、服务熔断与限流降级、自动化故障转移以及混沌工程故障主动注入测试。
\end{itemize}

\subsubsection{示范性学习成果与能力目标}
\begin{enumerate}
    \item \textbf{[解释]} 深入辨析软件程序编写缺陷、操作系统运行时差错、底层物理硬件故障以及运行时异常（Exceptions，如除零异常、越界中断）之间的层次传导关系与本质区别。
    \item \textbf{[解释]} 深刻阐释在系统工程中对故障进行“早感知检测”、“局部隔离容忍”以及“无缝状态自愈与恢复”的完整防御纵深与关键实现手段。
    \item \textbf{[解释]} 阐释纠错编码技术（ECC）在易失性动态内存（DRAM）、持久化存储介质以及网络物理信道通信中防范比特翻转衰减与噪声干扰的关键保护作用。
    \item \textbf{[应用]} 运用经典信息冗余算法（如海明编码、奇偶校验），设计并编程实现具备数据位差错检测与自修复能力的数据编解码软件模块。
    \item \textbf{[评估]} 从数据存储空间冗余开销、算法计算复杂度，以及编码、检错与纠错在不同系统层级上的时间延迟代价等维度，多角度系统对比不同纠错机制的综合优劣与工程选型策略。
\end{enumerate}

\subsection{系统安全}
\textbf{【学时建议】}：2 计算机核心学时，1 领域核心学时。

本单元研究计算机系统面临的软硬件安全威胁与攻击向量，阐述基于密码学底座、最小特权架构设计与全方位纵深防御的现代系统防护工程。

\subsubsection{核心主题}
\begin{itemize}
    \item \textbf{常见系统安全威胁与攻击向量}：计算机病毒、木马与蠕虫；分布式拒绝服务攻击（DDoS）；内存缓冲区溢出与特权提升漏洞；恶意窃听与中间人攻击（MITM）。
    \item \textbf{系统级安全防御对策体系}：
    \begin{itemize}
        \item 现代密码学基石：对称与非对称数据加密算法、数字证书与 PKI 信任链、不可逆单向哈希函数与数字签名技术。
        \item 安全体系结构设计：最小特权原则、可信计算基（TCB）、强制与自主访问控制（MAC/DAC）、硬件隔离安全区（如 Intel SGX、ARM TrustZone）。
        \item 动态防御与安全审计：入侵检测与防御系统（IDS/IPS）、下一代防火墙（NGFW）、安全审计与系统日志不可篡改追踪。
    \end{itemize}
\end{itemize}

\subsubsection{示范性学习成果与能力目标}
\begin{enumerate}
    \item \textbf{[解释]} 详细阐明计算系统面临的常见安全威胁及其攻击机理，并能结合业界真实网络空间安全攻击案例进行深入剖析与安全推演。
    \item \textbf{[解释/评估]} 全面阐述应对各类系统安全漏洞与威胁的主流系统级防护对策，深刻理解密码学加密、沙箱隔离、可信执行环境与访问控制列表在全栈安全加固中的协同机制与部署架构。
\end{enumerate}

\subsection{系统设计}
\textbf{【学时建议】}：2 计算机核心学时，1 领域核心学时。

本单元培养学生的宏观系统架构设计与工程权衡思维，确立高可用、高扩展、强健壮与安全的通用设计准则，并深度剖析业界经典工业级系统架构。

\subsubsection{核心主题}
\begin{itemize}
    \item \textbf{系统设计核心通用准则}：
    \begin{itemize}
        \item 活跃性：系统保证终将发生期望进展，不发生死锁、活锁或无限饥饿。
        \item 安全性：系统保证绝不发生未定义的非法破坏状态，数据不变性与完整性始终受控。
        \item 鲁棒性/健壮性：在面对畸形输入、部分节点瘫痪或硬件瞬态失效时保持优雅降级。
        \item 可扩展性：系统算力随资源横向堆叠或纵向升级呈现近线性扩容。
        \item 系统安全性：防御未授权访问、恶意篡改与隐私数据泄露。
    \end{itemize}
    \item \textbf{业界代表性工业级系统架构全景剖析}：
    \begin{itemize}
        \item Apache HTTP Server：多进程 Prefork 与多线程 Worker 高性能并发连接服务模型。
        \item Apache Spark：基于内存 RDD（弹性分布式数据集）的有向无环图（DAG）分布式任务调度与执行引擎。
        \item Linux 操作系统：基于虚拟文件系统（VFS）、统一设备抽象模型与模块化单内核的工业级稳健架构设计。
    \end{itemize}
\end{itemize}

\subsubsection{示范性学习成果与能力目标}
\begin{enumerate}
    \item \textbf{[解释]} 深入阐明现代计算系统设计中的通用评价准则（包括活跃性、安全性、鲁棒性、高可扩展性与系统安全性）的形式化内涵与工程度量。
    \item \textbf{[开发]} 给定特定应用系统的业务功能性需求与关键非功能性服务等级协议（SLO/SLA 约束），设计并输出高水准、组件接口清晰的顶层系统宏观架构设计方案。
    \item \textbf{[解释/评估]} 深入剖析并阐述业界代表性复杂软件系统（如 Linux 内核、Spark 分布式计算框架、高性能 Web 服务器）的底层体系架构、核心关键决策与其背后的工程权衡取舍。
\end{enumerate}

\section{专业素养与思维倾向}

\begin{itemize}
    \item \textbf{工程严谨性}：在学习、剖析、量化评估与调试大型计算系统时，必须始终对微观软硬件细节、跨层接口边界、时序竞态条件以及潜在的故障诱因保持高度的敏锐度与严密性，养成“差之毫厘，失之千里”的严谨工程品格。
    \item \textbf{动态适应性与权衡思维}：在架构设计与技术选型中展现出高度的灵活性与辩证思维。深刻认知世上不存在放之四海而皆准的“银弹”系统设计；面对迥异的业务需求、硬件约束与工作负载特征，能够审时度势，在吞吐量、响应时延、一致性、开发复杂度与成本之间做出最科学合理的工程设计决策。
\end{itemize}

\section{数学基础要求}

\subsection{必修数学基础}
\begin{itemize}
    \item \textbf{离散数学}：
    \begin{itemize}
        \item 集合论与二元关系（等价关系、偏序关系）。
        \item 基础图论（树结构、有向无环图 DAG、拓扑排序、图的连通性分析）。
        \item 数理逻辑基础（布尔代数运算、命题逻辑、真值表推导）。
    \end{itemize}
    \item \textbf{线性代数}：
    \begin{itemize}
        \item 基础矩阵与向量运算（矩阵乘法、状态转移矩阵、基变换与空间映射）。
    \end{itemize}
    \item \textbf{概率与数理统计}：
    \begin{itemize}
        \item 随机变量及其分布特性。
        \item 贝叶斯定理（Bayes' Theorem）与条件概率。
        \item 数学期望与方差/标准差。
        \item 累积分布函数（CDF）与概率密度函数（PDF），用于延迟分位数（P95/P99/P99.9）评测分析。
    \end{itemize}
\end{itemize}

\subsection{进阶建议数学基础}
\begin{itemize}
    \item \textbf{排队论基础}：M/M/1 排队模型、Little 定律、服务到达率与服务速率平衡，用于系统服务延迟与队列积压分析。
    \item \textbf{随机过程初步}：离散时间马尔可夫链（DTMC）与连续时间马尔可夫链（CTMC）状态转移平衡方程。
\end{itemize}

\section{跨领域共享概念与系统思维}

本知识领域作为多门骨干系统课程的公约数底座，与以下专业核心知识领域深度交织：
\begin{itemize}
    \item \textbf{网络与通信 (NC)}：协议分层模型、套接字 API 抽象、TCP/UDP 传输状态机、端到端可靠性与拥塞控制。
    \item \textbf{操作系统 (OS)}：进程与线程并发模型、虚拟内存分页管理、多级调度器设计、系统调用与特权级转换。
    \item \textbf{计算机体系结构与组织 (AR)}：微架构流水线、多级层次缓存、指令级并行、存储器一致性与软硬件接口边界。
    \item \textbf{并行与分布式计算 (PDC)}：多核并发控制、无锁同步、分布式状态机复制、强扩展与弱扩展评测。
\end{itemize}

\textbf{系统思维跨学科贯通核心维度：}
\begin{itemize}
    \item \textbf{软硬件协同设计}：打通算法抽象、编译运行时与底层微架构，破除单层优化瓶颈。
    \item \textbf{分层与抽象屏蔽}：通过严密的 API 与契约屏蔽复杂性，实现关注点分离。
    \item \textbf{状态与状态机模型}：统一利用有限状态机刻画时序控制流、通信交互协议与机器执行轨迹。
    \item \textbf{资源虚拟化与隔离}：以间接层提供透明独占抽象，通过沙箱机制实现多租户强安全与 QoS 隔离。
    \item \textbf{性能瓶颈与局部性}：洞悉光速物理极限与层级时延鸿沟，深度利用时空局部性与缓存体系攻克“访存墙”。
    \item \textbf{容错与冗余设计}：承认硬件物理与软件逻辑故障的必然性，基于校验、重试、信息编码与副本机制构筑高可用防御纵深。
    \item \textbf{全生命周期系统安全性}：遵循最小特权与深度防御原则，融合密码学机制构筑系统可信根基。
\end{itemize}

\section{真实能力任务规范}

\begin{taskbox}[任务 1：阐述系统设计的通用设计准则]
\textbf{【任务情境】}：针对企业拟建设的特定高并发业务系统（如大规模 C/S 客户端-服务器架构应用或分布式数据库集群），召开架构方案技术评审会。\\
\textbf{【能力陈述】}：面对给定的业务应用场景，系统性阐述并论证在设计此类计算系统时必须重点考量的核心通用准则（包括活跃性、安全性、鲁棒性、可扩展性与系统安全性），阐明各设计准则间的冲突与工程权衡取舍，形成严谨的设计规范书。\\
\textbf{【能力领域】}：系统 \quad \textbf{【能力单元】}：需求分析 / 架构设计 / 技术文档规范\\
\textbf{【关联知识单元】}：系统基础/系统设计、系统基础/系统可靠性、系统基础/系统安全\\
\textbf{【核心要求】}：计算机核心 \quad \textbf{【达标等级】}：\textbf{解释}
\end{taskbox}

\begin{taskbox}[任务 2：剖析与对比不同资源调度算法机制]
\textbf{【任务情境】}：操作系统内核或分布式容器调度集群需要为多租户多任务高并发场景合理分配关键计算资源（CPU、内存、网络带宽）。\\
\textbf{【能力陈述】}：给定主流调度算法族（如静态/动态优先级调度、最早截止时间优先 EDF、抢占式轮转 RR 与非抢占式调度等），深入剖析各算法的底层决策逻辑与调度时序，并从吞吐率、响应延迟、开销及抗饥饿能力等维度全面评估其各自的优缺点与适用场景。\\
\textbf{【能力领域】}：系统 / 计算理论 \quad \textbf{【能力单元】}：需求分析 / 技术文档 / 性能评估\\
\textbf{【关联知识单元】}：系统基础/资源分配与调度、系统基础/系统性能\\
\textbf{【核心要求】}：计算机核心 / 领域核心 \quad \textbf{【达标等级】}：\textbf{解释与评估}
\end{taskbox}

\begin{taskbox}[任务 3：面向特定应用场景选择并实现调度算法]
\textbf{【任务情境】}：针对具有特定延迟敏感性约束的特定业务环境（如实时音视频流转码分发系统、高并发 API 聚合网关或批处理数据清洗平台），默认通用调度策略无法满足特定 SLA/SLO 约束。\\
\textbf{【能力陈述】}：深入分析目标应用场景的负载特征与性能目标，精准选型或创新设计针对性的定制化调度算法，编程实现该调度引擎原型，并利用合成工作负载对其吞吐量、平均周转时间及尾延迟进行量化评测与调优。\\
\textbf{【能力领域】}：系统 / 应用 \quad \textbf{【能力单元】}：需求分析 / 架构设计 / 核心开发 / 质量测试 / 技术文档 / 综合评估\\
\textbf{【关联知识单元】}：系统基础/资源分配与调度、系统基础/系统性能、系统基础/系统性能评估、系统基础/系统设计\\
\textbf{【核心要求】}：领域核心 \quad \textbf{【达标等级】}：\textbf{应用、评估与开发}
\end{taskbox}

\begin{taskbox}[任务 4：端到端高可用系统架构设计与工程研发]
\textbf{【任务情境】}：从零研发一套满足生产环境严苛工业标准的高性能、高可用分布式服务系统（如轻量级分布式键值数据库或高并发异步 Web 服务器）。\\
\textbf{【能力陈述】}：面对给定的业务场景需求及多维度非功能性质量约束（高吞吐量、低延迟、99.99\% 高可用性容灾与多层级安全防护），主导自顶向下的高内聚低耦合系统架构设计，端到端编码实现核心功能模块，实施全套单元测试与集成压测，顺利完成系统交付验收。\\
\textbf{【能力领域】}：系统 / 应用 \quad \textbf{【能力单元】}：需求工程 / 架构设计 / 核心开发 / 系统测试 / 生产部署 / 模块集成 / 技术文档 / 性能评测 / 验收交付\\
\textbf{【关联知识单元】}：系统基础/系统基础概念、系统基础/系统设计、系统基础/系统可靠性、系统基础/系统安全、系统基础/系统性能评估\\
\textbf{【核心要求】}：领域核心 \quad \textbf{【达标等级】}：\textbf{应用、评估与开发}
\end{taskbox}

\begin{taskbox}[任务 5：评测指定系统的多维综合性能指标]
\textbf{【任务情境】}：针对已上线的工业级计算系统或开源大数据处理组件，需客观评定其在极限负载下的处理能力与抗压边界。\\
\textbf{【能力陈述】}：针对目标被测系统，综合运用理论性能上限解析推演与实证基准测试实验，对其核心性能指标（计算正确性、并发吞吐量 QPS、平均响应时延以及 P95/P99/P99.9 尾延迟等）开展严密的受控测量，并输出详实的量化评估分析报告。\\
\textbf{【能力领域】}：系统 / 计算理论 \quad \textbf{【能力单元】}：质量测试 / 部署调优 / 实验文档 / 性能评估\\
\textbf{【关联知识单元】}：系统基础/系统性能、系统基础/系统性能评估\\
\textbf{【核心要求】}：计算机核心 / 领域核心 \quad \textbf{【达标等级】}：\textbf{应用与评估}
\end{taskbox}

\begin{taskbox}[任务 6：精准诊断与定位复杂系统的性能瓶颈]
\textbf{【任务情境】}：处于生产环境中的大型分布式计算系统遭遇突发性性能恶化或系统抖动，扩容算力无法换来线性的吞吐增长。\\
\textbf{【能力陈述】}：针对在真实部署环境中运行的复杂系统，穿透全栈调用链路，熟练运用系统级性能剖析工具与理论分析推导，精准排查并断定限制系统吞吐与时延的核心瓶颈（如内存访存墙、CPU 算力耗尽、网络/磁盘 I/O 拥塞或锁争用开销）。\\
\textbf{【能力领域】}：系统 / 计算理论 \quad \textbf{【能力单元】}：架构设计 / 系统测试 / 环境部署 / 技术文档 / 性能评估\\
\textbf{【关联知识单元】}：系统基础/系统性能、系统基础/系统性能评估、系统基础/系统设计、系统基础/计算机系统概览\\
\textbf{【核心要求】}：领域核心 \quad \textbf{【达标等级】}：\textbf{应用、评估与开发}
\end{taskbox}

\begin{taskbox}[任务 7：基于深度工具链排查并修复分层系统软件缺陷]
\textbf{【任务情境】}：在具备复杂分层架构的软件系统中，偶发深层次跨层异常、并发竞态死锁或未知内存越界破坏，导致服务间歇性崩溃。\\
\textbf{【能力陈述】}：基于目标系统的故障异常日志与运行环境，熟练运用高级程序跟踪分析、单步调试、核心转储 Core Dump 与内存检测工具链，在庞大的调用栈中迅速定位深层次代码缺陷根源，并实施高质量的代码重构与回归验证测试。\\
\textbf{【能力领域】}：软件 / 系统 / 应用 \quad \textbf{【能力单元】}：核心开发 / 质量测试 / 文档记录 / 缺陷评估 / 运维保障\\
\textbf{【关联知识单元】}：系统基础/系统设计、系统基础/系统可靠性、系统基础/系统性能评估\\
\textbf{【核心要求】}：计算机核心 / 领域核心 \quad \textbf{【达标等级】}：\textbf{应用、评估与开发}
\end{taskbox}

\begin{taskbox}[任务 8：将本地系统平滑部署至现代云原生环境]
\textbf{【任务情境】}：研发团队需将原本在本地物理单机上开发测试完备的业务系统，向弹性公共云或私有云基础设施进行无缝云原生化迁移与多节点部署。\\
\textbf{【能力陈述】}：给定可在本地机器正常运行的系统程序，进行容器化封装镜像构建，编排配置容器网络拓扑、持久化数据卷与资源限额，顺利将其部署至基于 Kubernetes 或云平台容器编排引擎的集群环境中，达成高可用弹性运行。\\
\textbf{【能力领域】}：软件 / 系统 \quad \textbf{【能力单元】}：生产部署 / 模块集成 / 技术文档规范\\
\textbf{【关联知识单元】}：系统基础/计算机系统概览\\
\textbf{【核心要求】}：计算机核心 \quad \textbf{【达标等级】}：\textbf{应用}
\end{taskbox}

\begin{taskbox}[任务 9：大型系统的日常高可用运维与架构持续演进]
\textbf{【任务情境】}：面对工业级在线系统的生产常态化运行，业务流量急剧攀升，用户新功能需求与故障反馈纷至沓来，系统面临可用性与演进压力。\\
\textbf{【能力陈述】}：针对长期上线运转的大型系统，负责其全天候运维监控、告警响应、故障自愈与容灾灾备；并以用户反馈与生产监控数据为驱动，规划制定系统架构重构与平滑升级方案，推动复杂系统持续健壮演进。\\
\textbf{【能力领域】}：软件 / 系统 \quad \textbf{【能力单元】}：需求演进 / 架构重构 / 软件开发 / 质量测试 / 自动化部署 / 模块集成 / 文档演进 / 运行评估 / 生产运维 / 工程管理 / 用户验收 / 社会议题调适 / 持续改进\\
\textbf{【关联知识单元】}：系统基础/计算机系统概览、系统基础/系统基础概念、系统基础/系统设计、系统基础/系统可靠性、系统基础/系统安全、系统基础/系统性能评估\\
\textbf{【核心要求】}：领域核心 \quad \textbf{【达标等级】}：\textbf{应用、评估与开发}
\end{taskbox}

\section{课程设计与教学实施建议}

\subsection{导论型课程：计算机系统导论 (40 学时)}
本门课程面向大一第二学期或大二第一学期的计算机科学专业本科生，旨在建立系统思维的宏观基石：
\begin{itemize}
    \item \textbf{学时规划分配}：
    \begin{itemize}
        \item 计算机系统概述 --- 2 学时
        \item 系统基本概念 --- 6 学时
        \item 资源分配与调度 --- 4 学时
        \item 系统性能 --- 6 学时
        \item 性能评估 --- 6 学时
        \item 系统可靠性 --- 4 学时
        \item 系统设计 --- 6 学时
        \item 系统安全 --- 6 学时
    \end{itemize}
    \item \textbf{先修课程要求}：
    \begin{itemize}
        \item 离散数学中的集合与关系、基础图论与基本布尔逻辑。
        \item 线性代数中的基础矩阵与向量运算。
        \item 概率与数理统计中的随机变量、贝叶斯定理、数学期望与方差、累积分布函数（CDF）与概率密度函数（PDF）。
    \end{itemize}
    \item \textbf{达成能力陈述}：
    完成本课程学习的学生应能够：(1) 深刻理解计算机系统的全栈核心概念；(2) 掌握在性能、可靠性与安全性维度设计计算系统的关键工程准则；(3) 熟练根据官方技术文档，在真实环境中成功部署并科学评测代表性复杂工业系统（如 MySQL 数据库与 Apache Spark 计算框架）；(4) 具备独立设计并编程实现简易计算系统的工程能力（如交互式命令行程序、高并发轻量级 Web 服务器、简易持久化存储系统）。
\end{itemize}

\subsection{进阶高级课程：高级计算机系统架构与设计 (40--45 学时)}
本门课程面向大三或大四本科生，侧重大型分布式系统设计、微架构协同调优与前沿工程实战：
\begin{itemize}
    \item \textbf{学时规划分配}：
    \begin{itemize}
        \item 系统设计深度专题 --- 10 学时
        \item 操作系统交叉单元 (OS 核心单元深化) --- 2 学时
        \item 网络通信交叉单元 (NC 核心单元深化) --- 2 学时
        \item 并行分布式计算交叉单元 (PDC 核心单元深化) --- 2 学时
        \item 体系结构交叉单元 (AR 核心单元深化) --- 2 学时
        \item 系统可靠性与容错机制 --- 10 学时
        \item 系统性能优化与评测进阶 --- 6 学时
        \item 系统安全与纵深防御架构 --- 6 学时
    \end{itemize}
    \item \textbf{先修课程要求}：
    \begin{itemize}
        \item 排队论基础与随机过程初步。
        \item 系统基础（SF）导论型先修课程（或系统核心课程组合）。
    \end{itemize}
    \item \textbf{达成能力陈述}：
    完成本门高级进阶课程的学生应能够：(1) 对计算机系统设计的核心体系架构原则建立深刻的本质认知；(2) 将这些核心原则深度映射至经典系统（如 Linux 操作系统内核、SQL 数据库引擎与 TCP/IP 协议栈）以及现代前沿系统（如 Hadoop/Spark 分布式计算框架、大规模分布式存储与云原生基础设施）的工业级设计中；(3) 具备独立设计、研发落地与系统调优复杂工业级计算机系统的核心工程实战能力（如从零构建微型文件系统、高并发异步 I/O 网络服务器）。
\end{itemize}

\section{指导委员会与工作组成员}

\begin{itemize}
    \item \textbf{工作组主席}：
    \begin{itemize}
        \item 向乔，厦门大学
    \end{itemize}
    \item \textbf{工作组成员}：
    \begin{itemize}
        \item Doug Lea，纽约州立大学奥斯威戈分校, 美国纽约
        \item Monica D. Anderson，阿拉巴马大学, 美国阿拉巴马
        \item Matthias Hauswirth，卢加诺大学, 瑞士卢加诺
        \item 翟恩南，阿里巴巴集团, 中国杭州
        \item 刘宇彤，上海交通大学, 中国上海
    \end{itemize}
    \item \textbf{重要贡献者}：
    \begin{itemize}
        \item Michael S. Kirkpatrick，詹姆斯麦迪逊大学, 美国弗吉尼亚
        \item 孔令和，上海交通大学, 中国上海
    \end{itemize}
\end{itemize}

\begin{strip}
\section{附录：核心主题与技能要求等级对照表}

\noindent\begin{tabularx}{\textwidth}{p{0.18\textwidth} Y c c c}
\toprule
\textbf{知识单元} & \textbf{核心知识点与具体主题} & \textbf{认知等级} & \textbf{分类} & \textbf{学时} \\
\midrule
\multirow{3}{*}{\makecell[l]{系统基础/计算机系统\\概述}} 
& 计算机基本组件（门、触发器、寄存器、总线；数据通路+控制+存储器）与体系结构历史演进 & 解释 & 计算机核心 & \multirow{3}{*}{3} \\
& 作为计算范式的硬件逻辑、布尔表达式、逻辑化简与与或式（SOP）；编程抽象、接口与库调用 & 解释 & 计算机核心 & \\
& 应用程序与操作系统服务边界、RPC；系统调用与特权级交互；流水线原理与强/弱扩展性 & 解释 & 计算机核心 & \\
\midrule
\multirow{3}{*}{\makecell[l]{系统基础/系统基本\\概念}} 
& 数字 vs 模拟系统、离散 vs 连续系统对比；基本逻辑门、时钟节拍与布尔代数化简 & 应用 & 计算机核心 & \multirow{3}{*}{4} \\
& 状态与状态转移生命周期；有限状态机模型（DFA/NFA）；组合逻辑与时序逻辑 & 应用 & 计算机核心 & \\
& 计算机与网络协议作为状态机抽象；串行单线程 vs 应用级并行处理（请求级、多线程多服务器、流水线） & 应用 & 计算机核心 & \\
\midrule
\multirow{2}{*}{\makecell[l]{系统基础/资源分配\\与调度}} 
& 异构计算机资源分类（CPU 份额、内存空间、磁盘存储 I/O、网络传输带宽） & 解释 & 计算机核心 & 1 \\
& 经典调度算法机制（FCFS、优先级调度、公平共享、抢占式/非抢占式、EDF）及其优缺点综合权衡 & 解释 & 领域核心 & 2 \\
\midrule
\multirow{3}{*}{\makecell[l]{系统基础/系统性能}} 
& 物理时延层次（光速极限与纳秒时钟、内存 vs 磁盘 vs 网络传输时延鸿沟） & 应用 & 计算机核心 & \multirow{3}{*}{\makecell[c]{2 \\ (CS:2 \\ KA:2)}} \\
& 缓存机制与时间/空间局部性；跨层缓存与缓存一致性（数据库、OS、分布式系统、微架构） & 应用 & 计算机核心 & \\
& 处理器层次存储体系与 AMAT 公式；虚拟化与隔离动因；间接层哲学与虚拟内存/虚拟机机制 & 应用 & 领域核心 & \\
\midrule
\multirow{3}{*}{\makecell[l]{系统基础/性能评估}} 
& 性能指标体系（吞吐量、延迟、P99 尾延迟、资源饱和度）；工作负载建模与基准测试集 & 评估 & 计算机核心 & \multirow{3}{*}{\makecell[c]{2 \\ (CS:2 \\ KA:2)}} \\
& CPI 性能方程；阿姆达尔定律数学推导与并行极限；数量级分析（大 O 记号）与快慢路径分离 & 评估 & 计算机核心 & \\
& 底层微架构事件损耗（流水线停顿、缓存缺失、缺页异常）；分层平台评测挑战与微基准测试陷阱 & 评估 & 领域核心 & \\
\midrule
\multirow{2}{*}{\makecell[l]{系统基础/系统可靠性}} 
& 软件缺陷、系统差错与物理故障的本质区别与级联传导；检错与重试机制 & 评估 & 计算机核心 & 2 \\
& 基于冗余编码（ECC、CRC、FEC）与副本镜像复制的高可用机制；心跳、看门狗、熔断与故障转移 & 评估 & 领域核心 & 1 \\
\midrule
\multirow{2}{*}{\makecell[l]{系统基础/系统安全}} 
& 常见系统安全威胁全景（病毒、蠕虫、DDoS 攻击、恶意窃听、特权提升漏洞） & 评估 & 计算机核心 & 2 \\
& 系统级安全防御对策（现代密码学算法、安全架构与最小特权原则、IDS/IPS 与防火墙防护） & 评估 & 领域核心 & 1 \\
\midrule
\multirow{2}{*}{\makecell[l]{系统基础/系统设计}} 
& 系统设计核心通用衡量准则（活跃性、安全性、鲁棒性、高可扩展性与系统安全性） & 开发 & 计算机核心 & 2 \\
& 业界工业级代表性系统架构全景深度剖析（Apache HTTP Server、Apache Spark、Linux 内核） & 开发 & 领域核心 & 1 \\
\bottomrule
\end{tabularx}
\end{strip}
